% Eksempelfil for å se Høst-temaene på LaTeX. Kan ikke kompileres alene.
\documentclass[a4paper,11pt]{article}
\usepackage[norsk]{babel}
\usepackage{amsmath,amssymb}
\usepackage[per-mode=symbol]{siunitx}
\usepackage{tikz,pgfplots}
\usepackage{minted}
\usetikzlibrary{arrows.meta,calc}
\usepackage[colorlinks=true]{hyperref}

\definecolor{darkred}{HTML}{781414}
\newcommand{\dd}{\mathrm{d}}
\renewcommand{\vec}[1]{\boldsymbol{#1}}
\DeclareMathOperator{\sgn}{sgn}
\NewDocumentCommand{\kraft}{m O{net}}{\vv{F}_{#2}}
\newtheorem{teorem}{Teorem}[section]

\title{Banen til en planet}
\author{Dag-Morten}

\begin{document}
\maketitle
\tableofcontents

\section{Keplers lover}\label{sec:kepler}

Keplers 2.\ lov sier at arealet planeten sveiper over per tidsenhet er det samme i hele banen. I vendepunktene står hastigheten $\vv{v}$ normalt på avstanden $r$ til sola, og vi får \emph{arealhastigheten} direkte, se likning~\eqref{eq:areal} og \cite[s.~12]{griffiths2017}.\footnote{Se også \url{https://www.uio.no} og \href{https://ndla.no}{NDLA}.}

\subsection*{Energibevaring}

\begin{equation}
    \frac{1}{2}\, r_0 v_0 \Delta t = \frac{1}{2}\, r_1 v_1 \Delta t
    \qquad\Longrightarrow\qquad
    v_1 = v_0\,\frac{r_0}{r_1} .
    \label{eq:areal}
\end{equation}

Energibevaring mellom de to vendepunktene gir
\begin{align*}
    \frac{1}{2} m v_0^2 - \frac{\gamma M m}{r_0} &= \frac{1}{2} m v_1^2 - \frac{\gamma M m}{r_1} , \\
    \frac{1}{2} v_0^2 \left(1 - \frac{r_0^2}{r_1^2}\right) &= \gamma M \left(\frac{1}{r_0} - \frac{1}{r_1}\right) \quad \text{etter at vi har satt inn $v_1$.}
\end{align*}

Gravitasjonskraften er $\vv{F}_1 = -\frac{\gamma M m}{r^2}\,\hat{\symbf{r}}$, og ved jordoverflata er $g = \qty{9,81}{m/s^2}$ og $\num{6,63e-34}$. Inline-matte kan også skrives \( \sin^2\theta + \cos^2\theta = 1 \), og display-matte
\[
    \int_0^\infty e^{-\alpha x^2} \,\dd x = \frac{1}{2}\sqrt{\frac{\pi}{\alpha}} , \qquad
    f(x) = \begin{cases} 1 & x > 0 \\ 0 & x \le 0 \end{cases}
\]

\begin{teorem}[Arealloven]\label{thm:areal}
Arealhastigheten $\frac{\dd A}{\dd t}$ er konstant når kraften er sentral.
\end{teorem}

\begin{enumerate}[label=\alph*)]
    \item finn uttrykket for $v_1$ fra arealloven,
    \item sett det inn i energilikningen, se avsnitt~\ref{sec:kepler},
    \item løs for $v_0$ og sett inn tall helt til slutt.
\end{enumerate}

\begin{table}[htbp]
    \centering
    \begin{tabular}{@{}l S[table-format=2.2] c@{}}
        \toprule
        Planet & {Avstand (\unit{AU})} & Måner \\
        \midrule
        Jorda & 1,00 & 1 \\
        Mars  & 1,52 & 2 \\
        \bottomrule
    \end{tabular}
    \caption{Avstand til sola.}\label{tab:planeter}
\end{table}

\begin{figure}[htbp]
    \centering
    \begin{tikzpicture}[>=Stealth, scale=1.2]
        \draw[->, thick] (0,0) -- (3,0) node[right] {$x$};
        \draw[darkred, dashed] (0,0) circle (1.5cm);
        \fill (1.5,0) circle (2pt) node[above right] {$P$};
        \coordinate (A) at ($(0,0)!0.5!(3,0)$);
    \end{tikzpicture}
    \begin{tikzpicture}
        \begin{axis}[xlabel={$t$ (\unit{s})}, ylabel={$v$}]
            \addplot[domain=0:5, samples=50] {9.81*x};
        \end{axis}
    \end{tikzpicture}
    \caption{Banen og fartsgrafen.}\label{fig:bane}
\end{figure}

\begin{minted}{python}
import numpy as np

def fart(t, g=9.81):
    """Farten etter tiden t."""
    return g * t  # m/s
\end{minted}

\begin{verbatim}
Dette står som det står: \frac{1}{2} $x$
\end{verbatim}

Kode i teksten: \verb|print("hei")| og \texttt{numpy}, \textbf{fet} og \textit{kursiv}, 50\,\% og 100~\% og et \% tegn. Linjeskift\\[2pt]
med luft \hspace{1em} og \vspace{2mm} mer.

\input{kapittel1}
\bibliographystyle{plain}
\bibliography{kilder}
\end{document}
